% =====================================================================
%  11  蒸留塔設計 — 左上=諸元表／右上=実寸比タワー図／下は空ける
%  source_memo: 05/06/07章＋trial #201。脱ブタン C3H8/C4H10 36段33m P19.5 全凝縮 α大／
%   脱エタン C2H6/C3H8 61段52m P7.9 部分凝縮 α3.9／C3 C3H6/C3H8 117段94m P16.7 全凝縮 α1.07(最難)。
%  方針：タイトル・共通設計枠組み・モデル名(SM/HYSYS)は載せない。表(左上)＋実寸比タワー(右上)。
%   下半分は意図的に空ける(後で別要素)。色は最難C3の赤のみ。丸括弧禁止。
% =====================================================================
\begin{frame}[t]

  \vspace{-0.7em}

  \begin{columns}[T,onlytextwidth]
    % ===== 左上：諸元表（SM/HYSYS は載せない／塔高は図側）=========
    \begin{column}{0.64\textwidth}
      {\scriptsize
       \setlength{\tabcolsep}{3pt}\renewcommand{\arraystretch}{1.04}%
       \begin{tabular}{@{}l c c c@{}}
         \toprule
          & 脱ブタン塔 & 脱エタン塔 & {\color{SpRed}C3 スプリッタ} \\
         \midrule
         軽／重キー & $\mathrm{C_3H_8}/\mathrm{C_4H_{10}}$ & $\mathrm{C_2H_6}/\mathrm{C_3H_8}$ & {\color{SpRed}$\mathrm{C_3H_6}/\mathrm{C_3H_8}$} \\
         塔頂留分 & $\mathrm{C_3H_8}$ & \begin{tabular}[t]{@{}c@{}}$\mathrm{H_2}\!\cdot\!\mathrm{CH_4}$\\[1pt]$\mathrm{C_2H_4}\!\cdot\!\mathrm{C_2H_6}$\end{tabular} & {\color{SpRed}$\mathrm{C_3H_6}$} \\
         比揮発度 $\alpha$ & 大 & $\approx 3.9$ & {\color{SpRed}\boldmath$\approx 1.07$} \\
         操作圧力 [bar] & $19.5$ & $7.9$ & {\color{SpRed}$16.7$} \\
         理論段数 & $36$ & $61$ & {\color{SpRed}\boldmath$117$} \\
         塔径 [m] & $4.0$ & $5.9$ & {\color{SpRed}$7.4$} \\
         還流比 & $2.0$ & $11.8$ & {\color{SpRed}$12.0$} \\
         フィード段 & $28$ & $31$ & {\color{SpRed}$96$} \\
         凝縮器 & 全凝縮 & 部分凝縮 & {\color{SpRed}全凝縮} \\
         リボイラ [MW] & $15.6$ & $44.2$ & {\color{SpRed}\boldmath$60.8$} \\
         \bottomrule
       \end{tabular}}
    \end{column}

    % ===== 右上：実寸比タワー図（名前＋塔高のみ）=================
    \begin{column}{0.36\textwidth}
      \centering
      \begin{tikzpicture}[
        x=1cm, y=1cm, >=Stealth, font=\scriptsize,
        name/.style={align=center, font=\scriptsize\bfseries},
        nameR/.style={align=center, font=\scriptsize\bfseries, text=SpRed},
      ]
        % 幅は塔径 D[m] に比例（4.0/5.9/7.4）、高さは塔高に比例（33/52/94m）
        % 脱ブタン D4.0 H33m
        \def\hB{1.34}\def\cB{0.55}\def\wB{0.33}
        \fill[pattern=horizontal lines, pattern color=HeaderBlue!55] (\cB-\wB,0) rectangle (\cB+\wB,\hB);
        \draw[HeaderBlue, line width=1pt] (\cB-\wB,0) rectangle (\cB+\wB,\hB);
        \node[name, anchor=south] at (\cB,\hB+0.06) {脱ブタン\\$33$ m};
        % 脱エタン D5.9 H52m
        \def\hE{2.11}\def\cE{1.9}\def\wE{0.49}
        \fill[pattern=horizontal lines, pattern color=HeaderBlue!55] (\cE-\wE,0) rectangle (\cE+\wE,\hE);
        \draw[HeaderBlue, line width=1pt] (\cE-\wE,0) rectangle (\cE+\wE,\hE);
        \node[name, anchor=south] at (\cE,\hE+0.06) {脱エタン\\$52$ m};
        % C3 D7.4 H94m（最難・赤）
        \def\hC{3.81}\def\cC{3.25}\def\wC{0.61}
        \fill[pattern=horizontal lines, pattern color=SpRed!50] (\cC-\wC,0) rectangle (\cC+\wC,\hC);
        \draw[SpRed, line width=1.2pt] (\cC-\wC,0) rectangle (\cC+\wC,\hC);
        \node[nameR, anchor=south] at (\cC,\hC+0.06) {C3 スプリッタ\\$94$ m};
        % 基準線＋注記
        \draw[black!35, line width=0.5pt] (\cB-0.5,0) -- (\cC+\wC+0.2,0);
      \end{tikzpicture}
    \end{column}
  \end{columns}

  \vspace{0.1em}

  % ===== 下：3塔の段プロファイル（組成・温度）全幅 =====
  \centering
  \includegraphics[width=0.97\linewidth]{dist_profiles}\par
  \vspace{0.12em}
  {\footnotesize\raggedright
   {\color{ArrowBlue}\textbf{脱ブタン}}・{\color{DipeOrange}\textbf{脱エタン}}：少ない段で組成・温度が急変化 ＝ 沸点差が大きく分離容易\par
   {\color{SpRed}\textbf{C3 スプリッタ}}：$117$ 段でも組成・温度ほぼ一定 ＝ 沸点差 $6$ K・$\alpha\!\approx\!1.07$ で分離困難\par}

\end{frame}
